回溯、枚举和剪枝深度题卡 按题型拆成章内大节，但每个大节都先从 LeetCode 案例切入，再进入分流、案例矩阵、案例推导和实验复盘。这样读者看到的是从具体题推到规律，而不是先读抽象方法再猜它能用在哪。

\topic{回溯、枚举与剪枝}
这一节先看 LeetCode 17 电话号码的字母组合、LeetCode 22 括号生成、LeetCode 39 组合总和 这几道 LeetCode 题，再整理同一题型的分流和推导。阅读顺序不是先背原理，而是先看案例的暴力瓶颈，再把重复出现的观察抽象成方法。

\subsection{原理和适用条件}

以 LeetCode 17 电话号码的字母组合、LeetCode 22 括号生成、LeetCode 39 组合总和 为主线：LeetCode 17 电话号码的字母组合：模型是“每个数字对应若干字母，答案是所有按位选择形成的字符串。”，优化抓手是“暴力就是完整笛卡尔积，优化重点是路径长度和下标同步推进，避免在每层重新拼接大量中间字符串。”；LeetCode 22 括号生成：模型是“合法括号串的前缀中右括号数量不能超过左括号数量。”，优化抓手是“状态保存已用左括号和右括号数量，只有满足约束的选择才继续递归。”；LeetCode 39 组合总和：模型是“候选数字可重复使用，答案不关心排列顺序。”，优化抓手是“排序后按起点枚举，递归时仍传当前下标表示可以继续使用同一数字。”。这些案例共同推出的规律是：先写搜索树：路径是什么、选择列表是什么、什么时候结束。然后用排序、剩余容量、同层去重或合法性预处理剪掉不可能分支。 方法名只是复盘后的标签，真正要掌握的是从题面约束、暴力失败、状态设计到边界反例的推导链。

\begin{principlebox}{图示：从 LeetCode 案例到 回溯、枚举与剪枝推导链}
\begin{tabular}{@{}p{0.18\linewidth}p{0.74\linewidth}@{}}
暴力基线 & 枚举候选、完整模拟或逐个检查，先保证覆盖全部可能答案。\\
瓶颈定位 & 慢点在于大量分支在很早就已经不可能成功，仍然被继续展开。重复元素还会产生等价排列或组合。\\
结构选择 & 剪枝来自容量、剩余长度、排序后去重、合法性预检查和提前终止。组合题关注起点推进，排列题关注元素是否使用，分割题关注当前片段是否合法。\\
不变量 & 路径数组保存已经做出的选择；起点下标保证组合不回头；同层去重保证同一个位置层级不会选择相同值；已使用数组保证排列不重复使用同一元素。\\
代码落地 & 回溯代码虽然天然递归，但工程中要意识到深度。题目规模较小时递归可读性好；若要遵守严格工程栈限制，可以改成显式栈状态机。\\
\end{tabular}
\end{principlebox}

\subsection{LeetCode 案例分流}

\begin{principlebox}{从 LeetCode 案例分流：回溯、枚举与剪枝}
\textbf{子集和组合。}
代表案例：LeetCode 78 子集、LeetCode 90 子集 II、LeetCode 77 组合、LeetCode 216 组合总和 III。先看这些题的题面条件：答案不关心顺序，元素按起点向后选择。由此推出的处理方式是：递增起点防止重复排列，同层去重防止重复值重复生成。风险点：同层去重和路径去重不是一回事。

\textbf{排列。}
代表案例：LeetCode 46 全排列、LeetCode 47 全排列 II。先看这些题的题面条件：答案关心顺序，每个位置选择一个未使用元素。由此推出的处理方式是：used 数组记录路径中已使用对象。风险点：重复元素排列要排序并跳过同层等价选择。

\textbf{分割和构造。}
代表案例：LeetCode 131 分割回文串、LeetCode 93 复原 IP 地址、LeetCode 140 单词拆分 II。先看这些题的题面条件：每一步选择一个前缀、片段或局部结构。由此推出的处理方式是：检查片段合法性后进入下一位置。风险点：输出规模可能指数级，复杂度要包含输出。

\textbf{棋盘和约束搜索。}
代表案例：LeetCode 51 N 皇后、LeetCode 79 单词搜索、LeetCode 212 单词搜索 II。先看这些题的题面条件：选择会影响行列、对角线或局部访问状态。由此推出的处理方式是：用集合或位掩码保存约束，进入退出时成对修改。风险点：剪枝只能删除不可能成功的分支。

\end{principlebox}

\subsection{案例矩阵}

本节案例覆盖 LeetCode 17 电话号码的字母组合、LeetCode 22 括号生成、LeetCode 39 组合总和、LeetCode 40 组合总和 II、LeetCode 46 全排列、LeetCode 47 全排列 II、LeetCode 51 N 皇后、LeetCode 78 子集、LeetCode 79 单词搜索、LeetCode 90 子集 II、LeetCode 93 复原 IP 地址、LeetCode 131 分割回文串、LeetCode 140 单词拆分 II、LeetCode 212 单词搜索 II、LeetCode 216 组合总和 III、LeetCode 679 24 点游戏、LeetCode 698 划分为 k 个相等的子集。阅读时不要按题号背答案，而要先判断题面落在哪个分流场景：查询目标是什么，输入约束给了什么单调性或等价关系，暴力慢在哪里，优化结构保存了什么状态。

\begin{principlebox}{案例矩阵}
\textbf{LeetCode 17 电话号码的字母组合。}
每个数字对应若干字母，答案是所有按位选择形成的字符串。单点风险：空输入、包含 7 或 9、输入中不应出现 0 和 1、输出规模指数增长。

\textbf{LeetCode 22 括号生成。}
合法括号串的前缀中右括号数量不能超过左括号数量。单点风险：n 为零或一、右括号提前超过左括号、输出规模是卡特兰数。

\textbf{LeetCode 39 组合总和。}
候选数字可重复使用，答案不关心排列顺序。单点风险：目标小于最小数、刚好等于、候选有重复时题意如何处理、输出规模。

\textbf{LeetCode 40 组合总和 II。}
每个候选只能使用一次，输入中可能存在重复值。单点风险：全重复、目标为零、同层去重写成跨层去重、答案顺序不要求。

\textbf{LeetCode 46 全排列。}
排列关心元素顺序，每个位置从所有未使用元素中选择一个。单点风险：空数组、单元素、输出规模为 n 阶乘、元素值唯一假设。

\textbf{LeetCode 47 全排列 II。}
重复元素会让不同下标选择生成相同排列。单点风险：全重复、部分重复、同层去重条件写反、输出顺序不要求。

\textbf{LeetCode 51 N 皇后。}
每一行放一个皇后，列和两条对角线不能冲突。单点风险：n 为一、n 为二或三无解、对角线编号、回退时状态恢复。

\textbf{LeetCode 78 子集。}
每个元素都有选或不选两种选择，答案是所有路径节点。单点风险：空数组、单元素、输出数量为二的 n 次方、顺序不要求。

\textbf{LeetCode 79 单词搜索。}
网格路径不能重复使用同一格，状态包含位置、匹配下标和访问标记。单点风险：单字符单词、路径需要回退、重复字母、单元格不能复用。

\textbf{LeetCode 90 子集 II。}
重复元素会让不同下标路径生成相同集合。单点风险：全重复、空集、重复值相邻、去重条件写成同层而不是全局。

\textbf{LeetCode 93 复原 IP 地址。}
字符串要切成四段，每段必须是 0 到 255 且无非法前导零。单点风险：前导零、数值超过 255、剩余字符太多或太少、刚好四段。

\textbf{LeetCode 131 分割回文串。}
每次选择一个从当前位置开始的回文前缀。单点风险：空串、单字符、全相同字符、重复分割路径。

\textbf{LeetCode 140 单词拆分 II。}
需要输出所有由字典单词拼出的句子，输出规模本身可能指数级。单点风险：无解后缀、重复单词、路径拼接成本、答案数量主导复杂度。

\textbf{LeetCode 212 单词搜索 II。}
多个单词在同一棋盘上搜索，公共前缀可以共享。单点风险：重复单词、同一格不能复用、前缀是完整单词、剪枝删除空子树。

\textbf{LeetCode 216 组合总和 III。}
从一到九中选 k 个数和为 n，每个数最多用一次。单点风险：n 太小、n 太大、k 为零、路径长度已经超过 k。

\textbf{LeetCode 679 24 点游戏。}
四个数字必须全部使用，括号表达式可以看成每次选两个数合并成一个新数。单点风险：减法和除法顺序、除零、浮点误差、重复中间状态。

\textbf{LeetCode 698 划分为 k 个相等的子集。}
每个元素要分配到 k 个桶中，所有桶目标和相同。单点风险：存在元素大于目标、重复元素、空桶对称、k 为一。

\end{principlebox}

\subsection{共性落点}

\begin{principlebox}{本组共性落点}
\textbf{慢版本。} 暴力回溯会枚举整个搜索树。它不是错误，而是必须先写清楚选择列表、路径状态和终止条件，再讨论剪枝。
\textbf{瓶颈。} 慢点在于大量分支在很早就已经不可能成功，仍然被继续展开。重复元素还会产生等价排列或组合。
\textbf{状态字段。} 路径、起点或使用标记、剩余目标、答案集合、剪枝条件；进入递归和退出递归必须成对恢复。
\textbf{不变量。} 路径数组保存已经做出的选择；起点下标保证组合不回头；同层去重保证同一个位置层级不会选择相同值；已使用数组保证排列不重复使用同一元素。
\textbf{实现检查。} 剪枝条件写在扩展前；同层去重不要误写成全局去重；保存答案时复制当前路径。
\textbf{C++ 落地。} 回溯代码虽然天然递归，但工程中要意识到深度。题目规模较小时递归可读性好；若要遵守严格工程栈限制，可以改成显式栈状态机。
\textbf{证明抓手。} 证明从搜索树覆盖开始：每个合法答案有且只有一条路径，剪枝只删掉不可能成功的分支。
\end{principlebox}

\subsection{案例推导}

\noindent\textbf{LeetCode 17 电话号码的字母组合。}

\textbf{题目说明。} LeetCode 17 电话号码的字母组合 的题面入口要先还原成具体任务：每个数字对应若干字母，答案是所有按位选择形成的字符串。

\textbf{样例入口。} 拿 digits=23 手算，2 对应 a/b/c，3 对应 d/e/f，第一层选 a 后第二层依次选 d/e/f，得到 ad、ae、af。

\textbf{暴力基线。} 慢版本就是完整搜索树：每层列出选择列表，路径到达终止条件时检查答案。它先保证覆盖所有可能，再把明显无效或重复的分支剪掉。放到本题就是：先按“每个数字对应若干字母，答案是所有按位选择形成的字符串”完整试慢版本；样例里已经能看到“规律是层数对应输入数字位置，路径保存已经选择的字母；空输入应返回空结果而不是包含空串”，所以优化前必须先能说清慢版本枚举了哪些候选。

\textbf{暴力过程。} 拿 digits=23 手算，2 对应 a/b/c，3 对应 d/e/f，第一层选 a 后第二层依次选 d/e/f，得到 ad、ae、af。

\textbf{重复工作。} 题面模型：每个数字对应若干字母，答案是所有按位选择形成的字符串。暴力版的慢点要落到本题：慢版本会围绕“每个数字对应若干字母，答案是所有按位选择形成的字符串”展开完整搜索树。样例规律是“规律是层数对应输入数字位置，路径保存已经选择的字母；空输入应返回空结果而不是包含空串”，说明一层递归必须对应一个明确决策，剪枝只能删除整棵不可能成功或重复的子树。后面的优化只消掉这类重复工作，不能改变答案集合。

\textbf{优化条件。} 本题的关键观察是：暴力就是完整笛卡尔积，优化重点是路径长度和下标同步推进，避免在每层重新拼接大量中间字符串。这句话必须能翻译成可维护状态；如果题目条件改变后观察不再成立，就不能继续套原来的写法。

\textbf{相邻状态。} 规律是层数对应输入数字位置，路径保存已经选择的字母；空输入应返回空结果而不是包含空串。把这个特例继续拆开看：一层递归对应一个明确决策，路径保存已经做出的选择，选择列表保存仍可能扩展的分支。剪枝前必须能说明被剪掉的整棵子树没有合法答案，而不是只因为它看起来不优。

\textbf{状态复用。} 把完整搜索树加上合法性检查、剩余资源剪枝和同层去重；路径状态进入和退出要成对恢复。本题落点是：规律是层数对应输入数字位置，路径保存已经选择的字母；空输入应返回空结果而不是包含空串。手算时只改被新输入影响的那一小部分，而不是回头重算完整候选。

\textbf{指针和字段。} 状态字段要能说清：路径、起点或使用标记、剩余目标、答案集合、剪枝条件；进入递归和退出递归必须成对恢复。手算时按这个协议跟踪：拿 digits=23 手算，2 对应 a/b/c，3 对应 d/e/f，第一层选 a 后第二层依次选 d/e/f，得到 ad、ae、af。然后按协议复查：手算时给每层标出选择列表和路径。剪枝发生前要指出被剪分支为什么已经不可能得到合法答案。

\textbf{代码落点。} 剪枝条件写在扩展前；同层去重不要误写成全局去重；保存答案时复制当前路径。关键代码行必须能逐句翻译成“旧状态怎样变成新状态”，否则就只是模板。

\textbf{正确性。} 证明从搜索树覆盖开始：每个合法答案有且只有一条路径，剪枝只删掉不可能成功的分支。

\textbf{边界实验。} 先用空输入、包含 7 或 9、输入中不应出现 0 和 1、输出规模指数增长做反例检查。实验时覆盖空输入、包含 7 或 9、输入中不应出现 0 和 1、输出规模指数增长，统计搜索节点数量和剪枝次数。重复元素题要打印同一层跳过哪些值，确认没有误删合法路径。

\textbf{复杂度变化。} 暴力搜索树的规模由输出或选择空间决定；剪枝不改变最坏指数级本质，但能删除不可能成功或重复的分支，复杂度要说明输出规模。

\textbf{迁移追问。} 如果题目改成“若字符映射来自键盘以外的字典，只需要替换候选表，搜索骨架不变”，先判断原状态是否仍然覆盖新答案集合，再决定是微调边界、改 value 语义，还是换算法。

\noindent\textbf{LeetCode 22 括号生成。}

\textbf{题目说明。} LeetCode 22 括号生成 的题面入口要先还原成具体任务：合法括号串的前缀中右括号数量不能超过左括号数量。

\textbf{样例入口。} 拿 n=2 手算，第一步只能放左括号，之后可以放左括号得到 (())，也可以在已有左括号数量大于右括号时放右括号得到 ()()。

\textbf{暴力基线。} 慢版本就是完整搜索树：每层列出选择列表，路径到达终止条件时检查答案。它先保证覆盖所有可能，再把明显无效或重复的分支剪掉。放到本题就是：先按“合法括号串的前缀中右括号数量不能超过左括号数量”完整试慢版本；样例里已经能看到“规律是状态包含已放左括号数和右括号数，右括号不能超过左括号，左括号不能超过 n”，所以优化前必须先能说清慢版本枚举了哪些候选。

\textbf{暴力过程。} 拿 n=2 手算，第一步只能放左括号，之后可以放左括号得到 (())，也可以在已有左括号数量大于右括号时放右括号得到 ()()。

\textbf{重复工作。} 题面模型：合法括号串的前缀中右括号数量不能超过左括号数量。暴力版的慢点要落到本题：慢版本会围绕“合法括号串的前缀中右括号数量不能超过左括号数量”展开完整搜索树。样例规律是“规律是状态包含已放左括号数和右括号数，右括号不能超过左括号，左括号不能超过 n”，说明一层递归必须对应一个明确决策，剪枝只能删除整棵不可能成功或重复的子树。后面的优化只消掉这类重复工作，不能改变答案集合。

\textbf{优化条件。} 本题的关键观察是：状态保存已用左括号和右括号数量，只有满足约束的选择才继续递归。这句话必须能翻译成可维护状态；如果题目条件改变后观察不再成立，就不能继续套原来的写法。

\textbf{相邻状态。} 规律是状态包含已放左括号数和右括号数，右括号不能超过左括号，左括号不能超过 n。把这个特例继续拆开看：一层递归对应一个明确决策，路径保存已经做出的选择，选择列表保存仍可能扩展的分支。剪枝前必须能说明被剪掉的整棵子树没有合法答案，而不是只因为它看起来不优。

\textbf{状态复用。} 把完整搜索树加上合法性检查、剩余资源剪枝和同层去重；路径状态进入和退出要成对恢复。本题落点是：规律是状态包含已放左括号数和右括号数，右括号不能超过左括号，左括号不能超过 n。手算时只改被新输入影响的那一小部分，而不是回头重算完整候选。

\textbf{指针和字段。} 状态字段要能说清：路径、起点或使用标记、剩余目标、答案集合、剪枝条件；进入递归和退出递归必须成对恢复。手算时按这个协议跟踪：拿 n=2 手算，第一步只能放左括号，之后可以放左括号得到 (())，也可以在已有左括号数量大于右括号时放右括号得到 ()()。然后按协议复查：手算时给每层标出选择列表和路径。剪枝发生前要指出被剪分支为什么已经不可能得到合法答案。

\textbf{代码落点。} 剪枝条件写在扩展前；同层去重不要误写成全局去重；保存答案时复制当前路径。关键代码行必须能逐句翻译成“旧状态怎样变成新状态”，否则就只是模板。

\textbf{正确性。} 证明从搜索树覆盖开始：每个合法答案有且只有一条路径，剪枝只删掉不可能成功的分支。

\textbf{边界实验。} 先用n 为零或一、右括号提前超过左括号、输出规模是卡特兰数做反例检查。实验时覆盖n 为零或一、右括号提前超过左括号、输出规模是卡特兰数，统计搜索节点数量和剪枝次数。重复元素题要打印同一层跳过哪些值，确认没有误删合法路径。

\textbf{复杂度变化。} 暴力搜索树的规模由输出或选择空间决定；剪枝不改变最坏指数级本质，但能删除不可能成功或重复的分支，复杂度要说明输出规模。

\textbf{迁移追问。} 如果题目改成“若有多种括号类型，状态还需要栈保存类型匹配关系”，先判断原状态是否仍然覆盖新答案集合，再决定是微调边界、改 value 语义，还是换算法。

\noindent\textbf{LeetCode 39 组合总和。}

\textbf{题目说明。} LeetCode 39 组合总和 的题面入口要先还原成具体任务：候选数字可重复使用，答案不关心排列顺序。

\textbf{样例入口。} 拿 candidates=[2, 3, 6, 7]、target=7 手算，选择 2 后仍可以继续选择 2，因为每个数可重复；路径 2, 2, 3 成功，7 单独也成功。

\textbf{暴力基线。} 慢版本就是完整搜索树：每层列出选择列表，路径到达终止条件时检查答案。它先保证覆盖所有可能，再把明显无效或重复的分支剪掉。放到本题就是：先按“候选数字可重复使用，答案不关心排列顺序”完整试慢版本；样例里已经能看到“规律是回溯起点不回退，但递归下一层仍传当前下标以允许重复使用”，所以优化前必须先能说清慢版本枚举了哪些候选。

\textbf{暴力过程。} 拿 candidates=[2, 3, 6, 7]、target=7 手算，选择 2 后仍可以继续选择 2，因为每个数可重复；路径 2, 2, 3 成功，7 单独也成功。

\textbf{重复工作。} 题面模型：候选数字可重复使用，答案不关心排列顺序。暴力版的慢点要落到本题：慢版本会围绕“候选数字可重复使用，答案不关心排列顺序”展开完整搜索树。样例规律是“规律是回溯起点不回退，但递归下一层仍传当前下标以允许重复使用”，说明一层递归必须对应一个明确决策，剪枝只能删除整棵不可能成功或重复的子树。后面的优化只消掉这类重复工作，不能改变答案集合。

\textbf{优化条件。} 本题的关键观察是：排序后按起点枚举，递归时仍传当前下标表示可以继续使用同一数字。这句话必须能翻译成可维护状态；如果题目条件改变后观察不再成立，就不能继续套原来的写法。

\textbf{相邻状态。} 规律是回溯起点不回退，但递归下一层仍传当前下标以允许重复使用。把这个特例继续拆开看：一层递归对应一个明确决策，路径保存已经做出的选择，选择列表保存仍可能扩展的分支。剪枝前必须能说明被剪掉的整棵子树没有合法答案，而不是只因为它看起来不优。

\textbf{状态复用。} 把完整搜索树加上合法性检查、剩余资源剪枝和同层去重；路径状态进入和退出要成对恢复。本题落点是：规律是回溯起点不回退，但递归下一层仍传当前下标以允许重复使用。手算时只改被新输入影响的那一小部分，而不是回头重算完整候选。

\textbf{指针和字段。} 状态字段要能说清：路径、起点或使用标记、剩余目标、答案集合、剪枝条件；进入递归和退出递归必须成对恢复。手算时按这个协议跟踪：拿 candidates=[2, 3, 6, 7]、target=7 手算，选择 2 后仍可以继续选择 2，因为每个数可重复；路径 2, 2, 3 成功，7 单独也成功。然后按协议复查：手算时给每层标出选择列表和路径。剪枝发生前要指出被剪分支为什么已经不可能得到合法答案。

\textbf{代码落点。} 剪枝条件写在扩展前；同层去重不要误写成全局去重；保存答案时复制当前路径。关键代码行必须能逐句翻译成“旧状态怎样变成新状态”，否则就只是模板。

\textbf{正确性。} 证明从搜索树覆盖开始：每个合法答案有且只有一条路径，剪枝只删掉不可能成功的分支。

\textbf{边界实验。} 先用目标小于最小数、刚好等于、候选有重复时题意如何处理、输出规模做反例检查。实验时覆盖目标小于最小数、刚好等于、候选有重复时题意如何处理、输出规模，统计搜索节点数量和剪枝次数。重复元素题要打印同一层跳过哪些值，确认没有误删合法路径。

\textbf{复杂度变化。} 暴力搜索树的规模由输出或选择空间决定；剪枝不改变最坏指数级本质，但能删除不可能成功或重复的分支，复杂度要说明输出规模。

\textbf{迁移追问。} 如果题目改成“组合总和 II 每个数只能用一次，起点推进和去重条件都不同”，先判断原状态是否仍然覆盖新答案集合，再决定是微调边界、改 value 语义，还是换算法。

\noindent\textbf{LeetCode 40 组合总和 II。}

\textbf{题目说明。} LeetCode 40 组合总和 II 的题面入口要先还原成具体任务：每个候选只能使用一次，输入中可能存在重复值。

\textbf{样例入口。} 拿 candidates=[10, 1, 2, 7, 6, 1, 5]、target=8 手算，排序后两个 1 在同一层只能选第一个作为起点，否则会生成重复组合；但选了第一个 1 后，下一层可以选第二个 1。

\textbf{暴力基线。} 慢版本就是完整搜索树：每层列出选择列表，路径到达终止条件时检查答案。它先保证覆盖所有可能，再把明显无效或重复的分支剪掉。放到本题就是：先按“每个候选只能使用一次，输入中可能存在重复值”完整试慢版本；样例里已经能看到“规律是每个数只能用一次，递归下一层传 i+1，并做同层去重”，所以优化前必须先能说清慢版本枚举了哪些候选。

\textbf{暴力过程。} 拿 candidates=[10, 1, 2, 7, 6, 1, 5]、target=8 手算，排序后两个 1 在同一层只能选第一个作为起点，否则会生成重复组合；但选了第一个 1 后，下一层可以选第二个 1。

\textbf{重复工作。} 题面模型：每个候选只能使用一次，输入中可能存在重复值。暴力版的慢点要落到本题：慢版本会围绕“每个候选只能使用一次，输入中可能存在重复值”展开完整搜索树。样例规律是“规律是每个数只能用一次，递归下一层传 i+1，并做同层去重”，说明一层递归必须对应一个明确决策，剪枝只能删除整棵不可能成功或重复的子树。后面的优化只消掉这类重复工作，不能改变答案集合。

\textbf{优化条件。} 本题的关键观察是：排序后递归下一层从 i 加一开始，同层遇到相同值要跳过。这句话必须能翻译成可维护状态；如果题目条件改变后观察不再成立，就不能继续套原来的写法。

\textbf{相邻状态。} 规律是每个数只能用一次，递归下一层传 i+1，并做同层去重。把这个特例继续拆开看：一层递归对应一个明确决策，路径保存已经做出的选择，选择列表保存仍可能扩展的分支。剪枝前必须能说明被剪掉的整棵子树没有合法答案，而不是只因为它看起来不优。

\textbf{状态复用。} 把完整搜索树加上合法性检查、剩余资源剪枝和同层去重；路径状态进入和退出要成对恢复。本题落点是：规律是每个数只能用一次，递归下一层传 i+1，并做同层去重。手算时只改被新输入影响的那一小部分，而不是回头重算完整候选。

\textbf{指针和字段。} 状态字段要能说清：路径、起点或使用标记、剩余目标、答案集合、剪枝条件；进入递归和退出递归必须成对恢复。手算时按这个协议跟踪：拿 candidates=[10, 1, 2, 7, 6, 1, 5]、target=8 手算，排序后两个 1 在同一层只能选第一个作为起点，否则会生成重复组合；但选了第一个 1 后，下一层可以选第二个 1。然后按协议复查：手算时给每层标出选择列表和路径。剪枝发生前要指出被剪分支为什么已经不可能得到合法答案。

\textbf{代码落点。} 剪枝条件写在扩展前；同层去重不要误写成全局去重；保存答案时复制当前路径。关键代码行必须能逐句翻译成“旧状态怎样变成新状态”，否则就只是模板。

\textbf{正确性。} 证明从搜索树覆盖开始：每个合法答案有且只有一条路径，剪枝只删掉不可能成功的分支。

\textbf{边界实验。} 先用全重复、目标为零、同层去重写成跨层去重、答案顺序不要求做反例检查。实验时覆盖全重复、目标为零、同层去重写成跨层去重、答案顺序不要求，统计搜索节点数量和剪枝次数。重复元素题要打印同一层跳过哪些值，确认没有误删合法路径。

\textbf{复杂度变化。} 暴力搜索树的规模由输出或选择空间决定；剪枝不改变最坏指数级本质，但能删除不可能成功或重复的分支，复杂度要说明输出规模。

\textbf{迁移追问。} 如果题目改成“子集 II 使用同样的同层去重，只是每个路径节点都可能成为答案”，先判断原状态是否仍然覆盖新答案集合，再决定是微调边界、改 value 语义，还是换算法。

\noindent\textbf{LeetCode 46 全排列。}

\textbf{题目说明。} LeetCode 46 全排列 的题面入口要先还原成具体任务：排列关心元素顺序，每个位置从所有未使用元素中选择一个。

\textbf{样例入口。} 拿 [1, 2, 3] 手算，第一位可选 1、2、3；选 1 后，第二位只能从未使用的 2、3 中选。

\textbf{暴力基线。} 暴力可以枚举所有长度为 n 的序列，再检查是否每个元素刚好出现一次。

\textbf{暴力过程。} 以 [1, 2, 3] 为例，第一位可选 1、2、3；若第一位选 1，第二位只能从 2、3 中选，第三位由剩余元素决定。

\textbf{重复工作。} 题面模型：排列关心元素顺序，每个位置从所有未使用元素中选择一个。暴力版的慢点要落到本题：浪费在生成了含重复使用元素的非法序列。排列的每一层只需要从未使用元素中选择。后面的优化只消掉这类重复工作，不能改变答案集合。

\textbf{优化条件。} 本题的关键观察是：路径长度达到 n 时产生答案，used 数组保证同一元素不重复使用。这句话必须能翻译成可维护状态；如果题目条件改变后观察不再成立，就不能继续套原来的写法。

\textbf{相邻状态。} 规律是路径长度等于 n 时得到一个排列，used 数组表示哪些下标已经在当前路径里。把这个特例继续拆开看：一层递归对应一个明确决策，路径保存已经做出的选择，选择列表保存仍可能扩展的分支。剪枝前必须能说明被剪掉的整棵子树没有合法答案，而不是只因为它看起来不优。

\textbf{状态复用。} 回溯路径 path 保存当前排列前缀，used[i] 表示 nums[i] 是否已在路径中。path 长度等于 n 时复制进答案。手算时只改被新输入影响的那一小部分，而不是回头重算完整候选。

\textbf{指针和字段。} 状态字段要能说清：depth 是当前要填的位置，path 是已选择序列，used 是可选状态，answers 保存所有完整排列。手算时按这个协议跟踪：[1, 2, 3] 中路径 [1] 后，used[1]=true；下一层只能选 2 或 3。回退时弹出 1 并恢复 used。

\textbf{代码落点。} 关键是进入分支前 path.push\_back(nums[i]); used[i]=true；退出分支时反向恢复。保存答案时必须复制 path。关键代码行必须能逐句翻译成“旧状态怎样变成新状态”，否则就只是模板。

\textbf{正确性。} 每个位置都枚举所有未使用元素，所以任意合法排列都有唯一生成路径；used 阻止同一元素在一条路径中重复出现。

\textbf{边界实验。} 先用空数组、单元素、输出规模为 n 阶乘、元素值唯一假设做反例检查。实验时覆盖空数组、单元素、输出规模为 n 阶乘、元素值唯一假设，统计搜索节点数量和剪枝次数。重复元素题要打印同一层跳过哪些值，确认没有误删合法路径。

\textbf{复杂度变化。} 输出规模为 n! 个排列，每个排列长度 n，时间 O(n*n!)，空间 O(n) 不计输出。

\textbf{迁移追问。} 如果题目改成“若有重复元素，需要排序并在同层跳过等价选择”，先判断原状态是否仍然覆盖新答案集合，再决定是微调边界、改 value 语义，还是换算法。

\noindent\textbf{LeetCode 47 全排列 II。}

\textbf{题目说明。} LeetCode 47 全排列 II 的题面入口要先还原成具体任务：重复元素会让不同下标选择生成相同排列。

\textbf{样例入口。} 拿 [1, 1, 2] 手算，第一位选第一个 1 和第二个 1 会生成同样分支，所以同层必须跳过后一个 1；但第一层选 1 后，第二层仍可选另一个 1。

\textbf{暴力基线。} 慢版本就是完整搜索树：每层列出选择列表，路径到达终止条件时检查答案。它先保证覆盖所有可能，再把明显无效或重复的分支剪掉。放到本题就是：先按“重复元素会让不同下标选择生成相同排列”完整试慢版本；样例里已经能看到“规律是排序后用 used 区分路径内使用，同层重复值只保留第一个未使用代表”，所以优化前必须先能说清慢版本枚举了哪些候选。

\textbf{暴力过程。} 拿 [1, 1, 2] 手算，第一位选第一个 1 和第二个 1 会生成同样分支，所以同层必须跳过后一个 1；但第一层选 1 后，第二层仍可选另一个 1。

\textbf{重复工作。} 题面模型：重复元素会让不同下标选择生成相同排列。暴力版的慢点要落到本题：慢版本会围绕“重复元素会让不同下标选择生成相同排列”展开完整搜索树。样例规律是“规律是排序后用 used 区分路径内使用，同层重复值只保留第一个未使用代表”，说明一层递归必须对应一个明确决策，剪枝只能删除整棵不可能成功或重复的子树。后面的优化只消掉这类重复工作，不能改变答案集合。

\textbf{优化条件。} 本题的关键观察是：排序后使用 used 数组，同层遇到相同值且前一个相同值未使用时跳过。这句话必须能翻译成可维护状态；如果题目条件改变后观察不再成立，就不能继续套原来的写法。

\textbf{相邻状态。} 规律是排序后用 used 区分路径内使用，同层重复值只保留第一个未使用代表。把这个特例继续拆开看：一层递归对应一个明确决策，路径保存已经做出的选择，选择列表保存仍可能扩展的分支。剪枝前必须能说明被剪掉的整棵子树没有合法答案，而不是只因为它看起来不优。

\textbf{状态复用。} 把完整搜索树加上合法性检查、剩余资源剪枝和同层去重；路径状态进入和退出要成对恢复。本题落点是：规律是排序后用 used 区分路径内使用，同层重复值只保留第一个未使用代表。手算时只改被新输入影响的那一小部分，而不是回头重算完整候选。

\textbf{指针和字段。} 状态字段要能说清：路径、起点或使用标记、剩余目标、答案集合、剪枝条件；进入递归和退出递归必须成对恢复。手算时按这个协议跟踪：拿 [1, 1, 2] 手算，第一位选第一个 1 和第二个 1 会生成同样分支，所以同层必须跳过后一个 1；但第一层选 1 后，第二层仍可选另一个 1。然后按协议复查：手算时给每层标出选择列表和路径。剪枝发生前要指出被剪分支为什么已经不可能得到合法答案。

\textbf{代码落点。} 剪枝条件写在扩展前；同层去重不要误写成全局去重；保存答案时复制当前路径。关键代码行必须能逐句翻译成“旧状态怎样变成新状态”，否则就只是模板。

\textbf{正确性。} 证明从搜索树覆盖开始：每个合法答案有且只有一条路径，剪枝只删掉不可能成功的分支。

\textbf{边界实验。} 先用全重复、部分重复、同层去重条件写反、输出顺序不要求做反例检查。实验时覆盖全重复、部分重复、同层去重条件写反、输出顺序不要求，统计搜索节点数量和剪枝次数。重复元素题要打印同一层跳过哪些值，确认没有误删合法路径。

\textbf{复杂度变化。} 暴力搜索树的规模由输出或选择空间决定；剪枝不改变最坏指数级本质，但能删除不可能成功或重复的分支，复杂度要说明输出规模。

\textbf{迁移追问。} 如果题目改成“子集 II 和组合总和 II 也是同层去重，但起点推进方式不同”，先判断原状态是否仍然覆盖新答案集合，再决定是微调边界、改 value 语义，还是换算法。

\noindent\textbf{LeetCode 51 N 皇后。}

\textbf{题目说明。} LeetCode 51 N 皇后 的题面入口要先还原成具体任务：每一行放一个皇后，列和两条对角线不能冲突。

\textbf{样例入口。} 拿 n=4 手算，第一行放某列后，下一行不能放同列、主对角线和副对角线冲突的位置。

\textbf{暴力基线。} 慢版本就是完整搜索树：每层列出选择列表，路径到达终止条件时检查答案。它先保证覆盖所有可能，再把明显无效或重复的分支剪掉。放到本题就是：先按“每一行放一个皇后，列和两条对角线不能冲突”完整试慢版本；样例里已经能看到“规律是按行递归，每行放一个皇后，三个集合分别记录列、row-col、row+col 是否被占用”，所以优化前必须先能说清慢版本枚举了哪些候选。

\textbf{暴力过程。} 拿 n=4 手算，第一行放某列后，下一行不能放同列、主对角线和副对角线冲突的位置。

\textbf{重复工作。} 题面模型：每一行放一个皇后，列和两条对角线不能冲突。暴力版的慢点要落到本题：慢版本会围绕“每一行放一个皇后，列和两条对角线不能冲突”展开完整搜索树。样例规律是“规律是按行递归，每行放一个皇后，三个集合分别记录列、row-col、row+col 是否被占用”，说明一层递归必须对应一个明确决策，剪枝只能删除整棵不可能成功或重复的子树。后面的优化只消掉这类重复工作，不能改变答案集合。

\textbf{优化条件。} 本题的关键观察是：逐行搜索，用列集合、主对角线集合和副对角线集合快速判断合法性。这句话必须能翻译成可维护状态；如果题目条件改变后观察不再成立，就不能继续套原来的写法。

\textbf{相邻状态。} 规律是按行递归，每行放一个皇后，三个集合分别记录列、row-col、row+col 是否被占用。把这个特例继续拆开看：一层递归对应一个明确决策，路径保存已经做出的选择，选择列表保存仍可能扩展的分支。剪枝前必须能说明被剪掉的整棵子树没有合法答案，而不是只因为它看起来不优。

\textbf{状态复用。} 把完整搜索树加上合法性检查、剩余资源剪枝和同层去重；路径状态进入和退出要成对恢复。本题落点是：规律是按行递归，每行放一个皇后，三个集合分别记录列、row-col、row+col 是否被占用。手算时只改被新输入影响的那一小部分，而不是回头重算完整候选。

\textbf{指针和字段。} 状态字段要能说清：路径、起点或使用标记、剩余目标、答案集合、剪枝条件；进入递归和退出递归必须成对恢复。手算时按这个协议跟踪：拿 n=4 手算，第一行放某列后，下一行不能放同列、主对角线和副对角线冲突的位置。然后按协议复查：手算时给每层标出选择列表和路径。剪枝发生前要指出被剪分支为什么已经不可能得到合法答案。

\textbf{代码落点。} 剪枝条件写在扩展前；同层去重不要误写成全局去重；保存答案时复制当前路径。关键代码行必须能逐句翻译成“旧状态怎样变成新状态”，否则就只是模板。

\textbf{正确性。} 证明从搜索树覆盖开始：每个合法答案有且只有一条路径，剪枝只删掉不可能成功的分支。

\textbf{边界实验。} 先用n 为一、n 为二或三无解、对角线编号、回退时状态恢复做反例检查。实验时覆盖n 为一、n 为二或三无解、对角线编号、回退时状态恢复，统计搜索节点数量和剪枝次数。重复元素题要打印同一层跳过哪些值，确认没有误删合法路径。

\textbf{复杂度变化。} 暴力搜索树的规模由输出或选择空间决定；剪枝不改变最坏指数级本质，但能删除不可能成功或重复的分支，复杂度要说明输出规模。

\textbf{迁移追问。} 如果题目改成“可用位掩码压缩三个约束集合，提高常数性能”，先判断原状态是否仍然覆盖新答案集合，再决定是微调边界、改 value 语义，还是换算法。

\noindent\textbf{LeetCode 78 子集。}

\textbf{题目说明。} LeetCode 78 子集 的题面入口要先还原成具体任务：每个元素都有选或不选两种选择，答案是所有路径节点。

\textbf{样例入口。} 拿 [1, 2] 画搜索树，第一层可以不选 1 或选 1；在选 1 的分支里再决定 2，得到 [1] 和 [1, 2]。

\textbf{暴力基线。} 暴力可以枚举长度从 0 到 n 的所有组合，并为每种长度单独搜索。

\textbf{暴力过程。} 以 [1, 2, 3] 为例，空集、[1]、[2]、[3]、[1, 2] 等都是答案；一个元素只有选和不选两种状态。

\textbf{重复工作。} 题面模型：每个元素都有选或不选两种选择，答案是所有路径节点。暴力版的慢点要落到本题：浪费在按长度重复组织搜索。子集问题的每个搜索树节点本身就是一个答案，不必等到固定长度。后面的优化只消掉这类重复工作，不能改变答案集合。

\textbf{优化条件。} 本题的关键观察是：回溯按起点向后枚举，避免不同顺序生成同一集合。这句话必须能翻译成可维护状态；如果题目条件改变后观察不再成立，就不能继续套原来的写法。

\textbf{相邻状态。} 规律是每个元素对应选或不选，或用起点向后枚举，路径上的每个节点都是一个子集；空数组也要输出空集。把这个特例继续拆开看：一层递归对应一个明确决策，路径保存已经做出的选择，选择列表保存仍可能扩展的分支。剪枝前必须能说明被剪掉的整棵子树没有合法答案，而不是只因为它看起来不优。

\textbf{状态复用。} 回溯时先把当前 path 加入答案，再从 start 到末尾选择下一个元素，递归时 start=i+1，保证元素顺序递增。手算时只改被新输入影响的那一小部分，而不是回头重算完整候选。

\textbf{指针和字段。} 状态字段要能说清：start 是下一次可选择的最小下标，path 是当前子集，answers 保存所有已访问路径节点。手算时按这个协议跟踪：从空集开始记录 []；选择 1 后记录 [1]；再选择 2 记录 [1, 2]；回退后选择 3 得到 [1, 3]。

\textbf{代码落点。} 关键是每进入一个节点就记录 path，而不是只在叶子记录。递归下一层用 i+1，避免 [1, 2] 和 [2, 1] 重复。关键代码行必须能逐句翻译成“旧状态怎样变成新状态”，否则就只是模板。

\textbf{正确性。} 任意子集按原数组下标升序排列后，对应搜索树中唯一一条递增路径；每个路径节点都被记录，所以不漏不重。

\textbf{边界实验。} 先用空数组、单元素、输出数量为二的 n 次方、顺序不要求做反例检查。实验时覆盖空数组、单元素、输出数量为二的 n 次方、顺序不要求，统计搜索节点数量和剪枝次数。重复元素题要打印同一层跳过哪些值，确认没有误删合法路径。

\textbf{复杂度变化。} 共有 \ensuremath{2^{n}} 个子集，复制路径总成本 O(\ensuremath{n \cdot 2^{n}})，递归栈 O(n)，输出空间另计。

\textbf{迁移追问。} 如果题目改成“也可用位掩码枚举所有子集，适合 n 较小的场景”，先判断原状态是否仍然覆盖新答案集合，再决定是微调边界、改 value 语义，还是换算法。

\noindent\textbf{LeetCode 79 单词搜索。}

\textbf{题目说明。} LeetCode 79 单词搜索 的题面入口要先还原成具体任务：网格路径不能重复使用同一格，状态包含位置、匹配下标和访问标记。

\textbf{样例入口。} 拿 board=[A B; C D]、word=AB 手算，从 A 出发只能走到相邻的 B，不能跳到 D；若搜索 ABA，第二个 A 不能复用起点格。

\textbf{暴力基线。} 慢版本就是完整搜索树：每层列出选择列表，路径到达终止条件时检查答案。它先保证覆盖所有可能，再把明显无效或重复的分支剪掉。放到本题就是：先按“网格路径不能重复使用同一格，状态包含位置、匹配下标和访问标记”完整试慢版本；样例里已经能看到“规律是状态包含位置、匹配下标和访问标记，递归退出时要恢复访问标记”，所以优化前必须先能说清慢版本枚举了哪些候选。

\textbf{暴力过程。} 拿 board=[A B; C D]、word=AB 手算，从 A 出发只能走到相邻的 B，不能跳到 D；若搜索 ABA，第二个 A 不能复用起点格。

\textbf{重复工作。} 题面模型：网格路径不能重复使用同一格，状态包含位置、匹配下标和访问标记。暴力版的慢点要落到本题：慢版本会围绕“网格路径不能重复使用同一格，状态包含位置、匹配下标和访问标记”展开完整搜索树。样例规律是“规律是状态包含位置、匹配下标和访问标记，递归退出时要恢复访问标记”，说明一层递归必须对应一个明确决策，剪枝只能删除整棵不可能成功或重复的子树。后面的优化只消掉这类重复工作，不能改变答案集合。

\textbf{优化条件。} 本题的关键观察是：剪枝包括首字符不匹配、剩余长度不足、字符频次不足和越界访问。这句话必须能翻译成可维护状态；如果题目条件改变后观察不再成立，就不能继续套原来的写法。

\textbf{相邻状态。} 规律是状态包含位置、匹配下标和访问标记，递归退出时要恢复访问标记。把这个特例继续拆开看：一层递归对应一个明确决策，路径保存已经做出的选择，选择列表保存仍可能扩展的分支。剪枝前必须能说明被剪掉的整棵子树没有合法答案，而不是只因为它看起来不优。

\textbf{状态复用。} 把完整搜索树加上合法性检查、剩余资源剪枝和同层去重；路径状态进入和退出要成对恢复。本题落点是：规律是状态包含位置、匹配下标和访问标记，递归退出时要恢复访问标记。手算时只改被新输入影响的那一小部分，而不是回头重算完整候选。

\textbf{指针和字段。} 状态字段要能说清：路径、起点或使用标记、剩余目标、答案集合、剪枝条件；进入递归和退出递归必须成对恢复。手算时按这个协议跟踪：拿 board=[A B; C D]、word=AB 手算，从 A 出发只能走到相邻的 B，不能跳到 D；若搜索 ABA，第二个 A 不能复用起点格。然后按协议复查：手算时给每层标出选择列表和路径。剪枝发生前要指出被剪分支为什么已经不可能得到合法答案。

\textbf{代码落点。} 剪枝条件写在扩展前；同层去重不要误写成全局去重；保存答案时复制当前路径。关键代码行必须能逐句翻译成“旧状态怎样变成新状态”，否则就只是模板。

\textbf{正确性。} 证明从搜索树覆盖开始：每个合法答案有且只有一条路径，剪枝只删掉不可能成功的分支。

\textbf{边界实验。} 先用单字符单词、路径需要回退、重复字母、单元格不能复用做反例检查。实验时覆盖单字符单词、路径需要回退、重复字母、单元格不能复用，统计搜索节点数量和剪枝次数。重复元素题要打印同一层跳过哪些值，确认没有误删合法路径。

\textbf{复杂度变化。} 暴力搜索树的规模由输出或选择空间决定；剪枝不改变最坏指数级本质，但能删除不可能成功或重复的分支，复杂度要说明输出规模。

\textbf{迁移追问。} 如果题目改成“若要查很多单词，应使用 Trie 合并搜索前缀”，先判断原状态是否仍然覆盖新答案集合，再决定是微调边界、改 value 语义，还是换算法。

\noindent\textbf{LeetCode 90 子集 II。}

\textbf{题目说明。} LeetCode 90 子集 II 的题面入口要先还原成具体任务：重复元素会让不同下标路径生成相同集合。

\textbf{样例入口。} 拿 [1, 2, 2] 手算，排序后第二个 2 如果在同一层再次作为起点，会生成和第一个 2 同样的子集；但在已经选择第一个 2 的下一层，第二个 2 又可以被选出 [2, 2]。

\textbf{暴力基线。} 慢版本就是完整搜索树：每层列出选择列表，路径到达终止条件时检查答案。它先保证覆盖所有可能，再把明显无效或重复的分支剪掉。放到本题就是：先按“重复元素会让不同下标路径生成相同集合”完整试慢版本；样例里已经能看到“规律是同层跳过重复值，不是全局禁止重复值”，所以优化前必须先能说清慢版本枚举了哪些候选。

\textbf{暴力过程。} 拿 [1, 2, 2] 手算，排序后第二个 2 如果在同一层再次作为起点，会生成和第一个 2 同样的子集；但在已经选择第一个 2 的下一层，第二个 2 又可以被选出 [2, 2]。

\textbf{重复工作。} 题面模型：重复元素会让不同下标路径生成相同集合。暴力版的慢点要落到本题：慢版本会围绕“重复元素会让不同下标路径生成相同集合”展开完整搜索树。样例规律是“规律是同层跳过重复值，不是全局禁止重复值”，说明一层递归必须对应一个明确决策，剪枝只能删除整棵不可能成功或重复的子树。后面的优化只消掉这类重复工作，不能改变答案集合。

\textbf{优化条件。} 本题的关键观察是：排序后在同一层跳过相同值，保留不同层使用重复值的可能。这句话必须能翻译成可维护状态；如果题目条件改变后观察不再成立，就不能继续套原来的写法。

\textbf{相邻状态。} 规律是同层跳过重复值，不是全局禁止重复值。把这个特例继续拆开看：一层递归对应一个明确决策，路径保存已经做出的选择，选择列表保存仍可能扩展的分支。剪枝前必须能说明被剪掉的整棵子树没有合法答案，而不是只因为它看起来不优。

\textbf{状态复用。} 把完整搜索树加上合法性检查、剩余资源剪枝和同层去重；路径状态进入和退出要成对恢复。本题落点是：规律是同层跳过重复值，不是全局禁止重复值。手算时只改被新输入影响的那一小部分，而不是回头重算完整候选。

\textbf{指针和字段。} 状态字段要能说清：路径、起点或使用标记、剩余目标、答案集合、剪枝条件；进入递归和退出递归必须成对恢复。手算时按这个协议跟踪：拿 [1, 2, 2] 手算，排序后第二个 2 如果在同一层再次作为起点，会生成和第一个 2 同样的子集；但在已经选择第一个 2 的下一层，第二个 2 又可以被选出 [2, 2]。然后按协议复查：手算时给每层标出选择列表和路径。剪枝发生前要指出被剪分支为什么已经不可能得到合法答案。

\textbf{代码落点。} 剪枝条件写在扩展前；同层去重不要误写成全局去重；保存答案时复制当前路径。关键代码行必须能逐句翻译成“旧状态怎样变成新状态”，否则就只是模板。

\textbf{正确性。} 证明从搜索树覆盖开始：每个合法答案有且只有一条路径，剪枝只删掉不可能成功的分支。

\textbf{边界实验。} 先用全重复、空集、重复值相邻、去重条件写成同层而不是全局做反例检查。实验时覆盖全重复、空集、重复值相邻、去重条件写成同层而不是全局，统计搜索节点数量和剪枝次数。重复元素题要打印同一层跳过哪些值，确认没有误删合法路径。

\textbf{复杂度变化。} 暴力搜索树的规模由输出或选择空间决定；剪枝不改变最坏指数级本质，但能删除不可能成功或重复的分支，复杂度要说明输出规模。

\textbf{迁移追问。} 如果题目改成“组合总和 II 使用同样的同层去重思想”，先判断原状态是否仍然覆盖新答案集合，再决定是微调边界、改 value 语义，还是换算法。

\noindent\textbf{LeetCode 93 复原 IP 地址。}

\textbf{题目说明。} LeetCode 93 复原 IP 地址 的题面入口要先还原成具体任务：字符串要切成四段，每段必须是 0 到 255 且无非法前导零。

\textbf{样例入口。} 拿 25525511135 手算，第一段可以是 255，但不能是 2552；段值必须在 0..255，且 01 这种前导零非法。

\textbf{暴力基线。} 慢版本就是完整搜索树：每层列出选择列表，路径到达终止条件时检查答案。它先保证覆盖所有可能，再把明显无效或重复的分支剪掉。放到本题就是：先按“字符串要切成四段，每段必须是 0 到 255 且无非法前导零”完整试慢版本；样例里已经能看到“规律是回溯选四段，每段长度 1 到 3，剩余字符数必须能填满剩余段数”，所以优化前必须先能说清慢版本枚举了哪些候选。

\textbf{暴力过程。} 拿 25525511135 手算，第一段可以是 255，但不能是 2552；段值必须在 0..255，且 01 这种前导零非法。

\textbf{重复工作。} 题面模型：字符串要切成四段，每段必须是 0 到 255 且无非法前导零。暴力版的慢点要落到本题：慢版本会围绕“字符串要切成四段，每段必须是 0 到 255 且无非法前导零”展开完整搜索树。样例规律是“规律是回溯选四段，每段长度 1 到 3，剩余字符数必须能填满剩余段数”，说明一层递归必须对应一个明确决策，剪枝只能删除整棵不可能成功或重复的子树。后面的优化只消掉这类重复工作，不能改变答案集合。

\textbf{优化条件。} 本题的关键观察是：按段数和剩余长度剪枝，每次尝试长度一到三的片段。这句话必须能翻译成可维护状态；如果题目条件改变后观察不再成立，就不能继续套原来的写法。

\textbf{相邻状态。} 规律是回溯选四段，每段长度 1 到 3，剩余字符数必须能填满剩余段数。把这个特例继续拆开看：一层递归对应一个明确决策，路径保存已经做出的选择，选择列表保存仍可能扩展的分支。剪枝前必须能说明被剪掉的整棵子树没有合法答案，而不是只因为它看起来不优。

\textbf{状态复用。} 把完整搜索树加上合法性检查、剩余资源剪枝和同层去重；路径状态进入和退出要成对恢复。本题落点是：规律是回溯选四段，每段长度 1 到 3，剩余字符数必须能填满剩余段数。手算时只改被新输入影响的那一小部分，而不是回头重算完整候选。

\textbf{指针和字段。} 状态字段要能说清：路径、起点或使用标记、剩余目标、答案集合、剪枝条件；进入递归和退出递归必须成对恢复。手算时按这个协议跟踪：拿 25525511135 手算，第一段可以是 255，但不能是 2552；段值必须在 0..255，且 01 这种前导零非法。然后按协议复查：手算时给每层标出选择列表和路径。剪枝发生前要指出被剪分支为什么已经不可能得到合法答案。

\textbf{代码落点。} 剪枝条件写在扩展前；同层去重不要误写成全局去重；保存答案时复制当前路径。关键代码行必须能逐句翻译成“旧状态怎样变成新状态”，否则就只是模板。

\textbf{正确性。} 证明从搜索树覆盖开始：每个合法答案有且只有一条路径，剪枝只删掉不可能成功的分支。

\textbf{边界实验。} 先用前导零、数值超过 255、剩余字符太多或太少、刚好四段做反例检查。实验时覆盖前导零、数值超过 255、剩余字符太多或太少、刚好四段，统计搜索节点数量和剪枝次数。重复元素题要打印同一层跳过哪些值，确认没有误删合法路径。

\textbf{复杂度变化。} 暴力搜索树的规模由输出或选择空间决定；剪枝不改变最坏指数级本质，但能删除不可能成功或重复的分支，复杂度要说明输出规模。

\textbf{迁移追问。} 如果题目改成“分割回文串也是前缀分割，但合法性检查不同”，先判断原状态是否仍然覆盖新答案集合，再决定是微调边界、改 value 语义，还是换算法。

\noindent\textbf{LeetCode 131 分割回文串。}

\textbf{题目说明。} LeetCode 131 分割回文串 的题面入口要先还原成具体任务：每次选择一个从当前位置开始的回文前缀。

\textbf{样例入口。} 拿 aab 手算，第一刀可以切 a，后面 ab 不全是回文；也可以切 aa，后面 b 是回文，得到 [aa, b]。

\textbf{暴力基线。} 慢版本就是完整搜索树：每层列出选择列表，路径到达终止条件时检查答案。它先保证覆盖所有可能，再把明显无效或重复的分支剪掉。放到本题就是：先按“每次选择一个从当前位置开始的回文前缀”完整试慢版本；样例里已经能看到“规律是回溯枚举切分点，但每段必须先通过回文检查；预处理回文表可以避免重复检查子串”，所以优化前必须先能说清慢版本枚举了哪些候选。

\textbf{暴力过程。} 拿 aab 手算，第一刀可以切 a，后面 ab 不全是回文；也可以切 aa，后面 b 是回文，得到 [aa, b]。

\textbf{重复工作。} 题面模型：每次选择一个从当前位置开始的回文前缀。暴力版的慢点要落到本题：慢版本会围绕“每次选择一个从当前位置开始的回文前缀”展开完整搜索树。样例规律是“规律是回溯枚举切分点，但每段必须先通过回文检查；预处理回文表可以避免重复检查子串”，说明一层递归必须对应一个明确决策，剪枝只能删除整棵不可能成功或重复的子树。后面的优化只消掉这类重复工作，不能改变答案集合。

\textbf{优化条件。} 本题的关键观察是：预处理回文 DP 表，让回溯合法性检查为常数。这句话必须能翻译成可维护状态；如果题目条件改变后观察不再成立，就不能继续套原来的写法。

\textbf{相邻状态。} 规律是回溯枚举切分点，但每段必须先通过回文检查；预处理回文表可以避免重复检查子串。把这个特例继续拆开看：一层递归对应一个明确决策，路径保存已经做出的选择，选择列表保存仍可能扩展的分支。剪枝前必须能说明被剪掉的整棵子树没有合法答案，而不是只因为它看起来不优。

\textbf{状态复用。} 把完整搜索树加上合法性检查、剩余资源剪枝和同层去重；路径状态进入和退出要成对恢复。本题落点是：规律是回溯枚举切分点，但每段必须先通过回文检查；预处理回文表可以避免重复检查子串。手算时只改被新输入影响的那一小部分，而不是回头重算完整候选。

\textbf{指针和字段。} 状态字段要能说清：路径、起点或使用标记、剩余目标、答案集合、剪枝条件；进入递归和退出递归必须成对恢复。手算时按这个协议跟踪：拿 aab 手算，第一刀可以切 a，后面 ab 不全是回文；也可以切 aa，后面 b 是回文，得到 [aa, b]。然后按协议复查：手算时给每层标出选择列表和路径。剪枝发生前要指出被剪分支为什么已经不可能得到合法答案。

\textbf{代码落点。} 剪枝条件写在扩展前；同层去重不要误写成全局去重；保存答案时复制当前路径。关键代码行必须能逐句翻译成“旧状态怎样变成新状态”，否则就只是模板。

\textbf{正确性。} 证明从搜索树覆盖开始：每个合法答案有且只有一条路径，剪枝只删掉不可能成功的分支。

\textbf{边界实验。} 先用空串、单字符、全相同字符、重复分割路径做反例检查。实验时覆盖空串、单字符、全相同字符、重复分割路径，统计搜索节点数量和剪枝次数。重复元素题要打印同一层跳过哪些值，确认没有误删合法路径。

\textbf{复杂度变化。} 暴力搜索树的规模由输出或选择空间决定；剪枝不改变最坏指数级本质，但能删除不可能成功或重复的分支，复杂度要说明输出规模。

\textbf{迁移追问。} 如果题目改成“若只求最少切割次数，可改成 DP 最短切分”，先判断原状态是否仍然覆盖新答案集合，再决定是微调边界、改 value 语义，还是换算法。

\noindent\textbf{LeetCode 140 单词拆分 II。}

\textbf{题目说明。} LeetCode 140 单词拆分 II 的题面入口要先还原成具体任务：需要输出所有由字典单词拼出的句子，输出规模本身可能指数级。

\textbf{样例入口。} 拿 catsanddog 手算，前缀 cat 和 cats 都能走，但某些后缀继续切会失败。若直接 DFS，每次都会重复探索同一个后缀。

\textbf{暴力基线。} 慢版本就是完整搜索树：每层列出选择列表，路径到达终止条件时检查答案。它先保证覆盖所有可能，再把明显无效或重复的分支剪掉。放到本题就是：先按“需要输出所有由字典单词拼出的句子，输出规模本身可能指数级”完整试慢版本；样例里已经能看到“规律是先用记忆化记录某个起点之后能生成哪些句子，或先用 DP 剪掉不可拆后缀，再枚举路径”，所以优化前必须先能说清慢版本枚举了哪些候选。

\textbf{暴力过程。} 拿 catsanddog 手算，前缀 cat 和 cats 都能走，但某些后缀继续切会失败。若直接 DFS，每次都会重复探索同一个后缀。

\textbf{重复工作。} 题面模型：需要输出所有由字典单词拼出的句子，输出规模本身可能指数级。暴力版的慢点要落到本题：慢版本会围绕“需要输出所有由字典单词拼出的句子，输出规模本身可能指数级”展开完整搜索树。样例规律是“规律是先用记忆化记录某个起点之后能生成哪些句子，或先用 DP 剪掉不可拆后缀，再枚举路径”，说明一层递归必须对应一个明确决策，剪枝只能删除整棵不可能成功或重复的子树。后面的优化只消掉这类重复工作，不能改变答案集合。

\textbf{优化条件。} 本题的关键观察是：先用 DP 或记忆化判断后缀是否可拆，再 DFS 枚举可行前缀并在末尾拼接句子。这句话必须能翻译成可维护状态；如果题目条件改变后观察不再成立，就不能继续套原来的写法。

\textbf{相邻状态。} 规律是先用记忆化记录某个起点之后能生成哪些句子，或先用 DP 剪掉不可拆后缀，再枚举路径。把这个特例继续拆开看：一层递归对应一个明确决策，路径保存已经做出的选择，选择列表保存仍可能扩展的分支。剪枝前必须能说明被剪掉的整棵子树没有合法答案，而不是只因为它看起来不优。

\textbf{状态复用。} 把完整搜索树加上合法性检查、剩余资源剪枝和同层去重；路径状态进入和退出要成对恢复。本题落点是：规律是先用记忆化记录某个起点之后能生成哪些句子，或先用 DP 剪掉不可拆后缀，再枚举路径。手算时只改被新输入影响的那一小部分，而不是回头重算完整候选。

\textbf{指针和字段。} 状态字段要能说清：路径、起点或使用标记、剩余目标、答案集合、剪枝条件；进入递归和退出递归必须成对恢复。手算时按这个协议跟踪：拿 catsanddog 手算，前缀 cat 和 cats 都能走，但某些后缀继续切会失败。若直接 DFS，每次都会重复探索同一个后缀。然后按协议复查：手算时给每层标出选择列表和路径。剪枝发生前要指出被剪分支为什么已经不可能得到合法答案。

\textbf{代码落点。} 剪枝条件写在扩展前；同层去重不要误写成全局去重；保存答案时复制当前路径。关键代码行必须能逐句翻译成“旧状态怎样变成新状态”，否则就只是模板。

\textbf{正确性。} 证明从搜索树覆盖开始：每个合法答案有且只有一条路径，剪枝只删掉不可能成功的分支。

\textbf{边界实验。} 先用无解后缀、重复单词、路径拼接成本、答案数量主导复杂度做反例检查。实验时覆盖无解后缀、重复单词、路径拼接成本、答案数量主导复杂度，统计搜索节点数量和剪枝次数。重复元素题要打印同一层跳过哪些值，确认没有误删合法路径。

\textbf{复杂度变化。} 暴力搜索树的规模由输出或选择空间决定；剪枝不改变最坏指数级本质，但能删除不可能成功或重复的分支，复杂度要说明输出规模。

\textbf{迁移追问。} 如果题目改成“只问能否拆分时用布尔 DP 即可，不要承担枚举所有句子的成本”，先判断原状态是否仍然覆盖新答案集合，再决定是微调边界、改 value 语义，还是换算法。

\noindent\textbf{LeetCode 212 单词搜索 II。}

\textbf{题目说明。} LeetCode 212 单词搜索 II 的题面入口要先还原成具体任务：多个单词在同一棋盘上搜索，公共前缀可以共享。

\textbf{样例入口。} 拿 board 上有 oath、pea、eat 手算，单词共享前缀时不应对每个词重复全图搜索；从 o 出发沿 Trie 前缀 o-a-t-h 找到 oath。

\textbf{暴力基线。} 慢版本就是完整搜索树：每层列出选择列表，路径到达终止条件时检查答案。它先保证覆盖所有可能，再把明显无效或重复的分支剪掉。放到本题就是：先按“多个单词在同一棋盘上搜索，公共前缀可以共享”完整试慢版本；样例里已经能看到“规律是 Trie 合并所有单词前缀，DFS 走棋盘同时走 Trie，命中单词后记录并可去重”，所以优化前必须先能说清慢版本枚举了哪些候选。

\textbf{暴力过程。} 拿 board 上有 oath、pea、eat 手算，单词共享前缀时不应对每个词重复全图搜索；从 o 出发沿 Trie 前缀 o-a-t-h 找到 oath。

\textbf{重复工作。} 题面模型：多个单词在同一棋盘上搜索，公共前缀可以共享。暴力版的慢点要落到本题：慢版本会围绕“多个单词在同一棋盘上搜索，公共前缀可以共享”展开完整搜索树。样例规律是“规律是 Trie 合并所有单词前缀，DFS 走棋盘同时走 Trie，命中单词后记录并可去重”，说明一层递归必须对应一个明确决策，剪枝只能删除整棵不可能成功或重复的子树。后面的优化只消掉这类重复工作，不能改变答案集合。

\textbf{优化条件。} 本题的关键观察是：先建 Trie，再从每个格子回溯；路径到达单词结束节点就收集答案。这句话必须能翻译成可维护状态；如果题目条件改变后观察不再成立，就不能继续套原来的写法。

\textbf{相邻状态。} 规律是 Trie 合并所有单词前缀，DFS 走棋盘同时走 Trie，命中单词后记录并可去重。把这个特例继续拆开看：一层递归对应一个明确决策，路径保存已经做出的选择，选择列表保存仍可能扩展的分支。剪枝前必须能说明被剪掉的整棵子树没有合法答案，而不是只因为它看起来不优。

\textbf{状态复用。} 把完整搜索树加上合法性检查、剩余资源剪枝和同层去重；路径状态进入和退出要成对恢复。本题落点是：规律是 Trie 合并所有单词前缀，DFS 走棋盘同时走 Trie，命中单词后记录并可去重。手算时只改被新输入影响的那一小部分，而不是回头重算完整候选。

\textbf{指针和字段。} 状态字段要能说清：路径、起点或使用标记、剩余目标、答案集合、剪枝条件；进入递归和退出递归必须成对恢复。手算时按这个协议跟踪：拿 board 上有 oath、pea、eat 手算，单词共享前缀时不应对每个词重复全图搜索；从 o 出发沿 Trie 前缀 o-a-t-h 找到 oath。然后按协议复查：手算时给每层标出选择列表和路径。剪枝发生前要指出被剪分支为什么已经不可能得到合法答案。

\textbf{代码落点。} 剪枝条件写在扩展前；同层去重不要误写成全局去重；保存答案时复制当前路径。关键代码行必须能逐句翻译成“旧状态怎样变成新状态”，否则就只是模板。

\textbf{正确性。} 证明从搜索树覆盖开始：每个合法答案有且只有一条路径，剪枝只删掉不可能成功的分支。

\textbf{边界实验。} 先用重复单词、同一格不能复用、前缀是完整单词、剪枝删除空子树做反例检查。实验时覆盖重复单词、同一格不能复用、前缀是完整单词、剪枝删除空子树，统计搜索节点数量和剪枝次数。重复元素题要打印同一层跳过哪些值，确认没有误删合法路径。

\textbf{复杂度变化。} 暴力搜索树的规模由输出或选择空间决定；剪枝不改变最坏指数级本质，但能删除不可能成功或重复的分支，复杂度要说明输出规模。

\textbf{迁移追问。} 如果题目改成“这是单词搜索从单目标扩展到多目标后的结构升级”，先判断原状态是否仍然覆盖新答案集合，再决定是微调边界、改 value 语义，还是换算法。

\noindent\textbf{LeetCode 216 组合总和 III。}

\textbf{题目说明。} LeetCode 216 组合总和 III 的题面入口要先还原成具体任务：从一到九中选 k 个数和为 n，每个数最多用一次。

\textbf{样例入口。} 拿 k=3、n=7 手算，选择 1 后还需两个数凑 6，可以选 2 和 4；若当前和已经超过 7，就不用再往下扩展。

\textbf{暴力基线。} 慢版本就是完整搜索树：每层列出选择列表，路径到达终止条件时检查答案。它先保证覆盖所有可能，再把明显无效或重复的分支剪掉。放到本题就是：先按“从一到九中选 k 个数和为 n，每个数最多用一次”完整试慢版本；样例里已经能看到“规律是路径长度、剩余和和起点共同定义状态，数字只能向后选避免重复组合”，所以优化前必须先能说清慢版本枚举了哪些候选。

\textbf{暴力过程。} 拿 k=3、n=7 手算，选择 1 后还需两个数凑 6，可以选 2 和 4；若当前和已经超过 7，就不用再往下扩展。

\textbf{重复工作。} 题面模型：从一到九中选 k 个数和为 n，每个数最多用一次。暴力版的慢点要落到本题：慢版本会围绕“从一到九中选 k 个数和为 n，每个数最多用一次”展开完整搜索树。样例规律是“规律是路径长度、剩余和和起点共同定义状态，数字只能向后选避免重复组合”，说明一层递归必须对应一个明确决策，剪枝只能删除整棵不可能成功或重复的子树。后面的优化只消掉这类重复工作，不能改变答案集合。

\textbf{优化条件。} 本题的关键观察是：按递增起点枚举，利用剩余个数和剩余和剪枝。这句话必须能翻译成可维护状态；如果题目条件改变后观察不再成立，就不能继续套原来的写法。

\textbf{相邻状态。} 规律是路径长度、剩余和和起点共同定义状态，数字只能向后选避免重复组合。把这个特例继续拆开看：一层递归对应一个明确决策，路径保存已经做出的选择，选择列表保存仍可能扩展的分支。剪枝前必须能说明被剪掉的整棵子树没有合法答案，而不是只因为它看起来不优。

\textbf{状态复用。} 把完整搜索树加上合法性检查、剩余资源剪枝和同层去重；路径状态进入和退出要成对恢复。本题落点是：规律是路径长度、剩余和和起点共同定义状态，数字只能向后选避免重复组合。手算时只改被新输入影响的那一小部分，而不是回头重算完整候选。

\textbf{指针和字段。} 状态字段要能说清：路径、起点或使用标记、剩余目标、答案集合、剪枝条件；进入递归和退出递归必须成对恢复。手算时按这个协议跟踪：拿 k=3、n=7 手算，选择 1 后还需两个数凑 6，可以选 2 和 4；若当前和已经超过 7，就不用再往下扩展。然后按协议复查：手算时给每层标出选择列表和路径。剪枝发生前要指出被剪分支为什么已经不可能得到合法答案。

\textbf{代码落点。} 剪枝条件写在扩展前；同层去重不要误写成全局去重；保存答案时复制当前路径。关键代码行必须能逐句翻译成“旧状态怎样变成新状态”，否则就只是模板。

\textbf{正确性。} 证明从搜索树覆盖开始：每个合法答案有且只有一条路径，剪枝只删掉不可能成功的分支。

\textbf{边界实验。} 先用n 太小、n 太大、k 为零、路径长度已经超过 k做反例检查。实验时覆盖n 太小、n 太大、k 为零、路径长度已经超过 k，统计搜索节点数量和剪枝次数。重复元素题要打印同一层跳过哪些值，确认没有误删合法路径。

\textbf{复杂度变化。} 暴力搜索树的规模由输出或选择空间决定；剪枝不改变最坏指数级本质，但能删除不可能成功或重复的分支，复杂度要说明输出规模。

\textbf{迁移追问。} 如果题目改成“可预估剩余最小和最大和进一步剪枝”，先判断原状态是否仍然覆盖新答案集合，再决定是微调边界、改 value 语义，还是换算法。

\noindent\textbf{LeetCode 679 24 点游戏。}

\textbf{题目说明。} LeetCode 679 24 点游戏 的题面入口要先还原成具体任务：四个数字必须全部使用，括号表达式可以看成每次选两个数合并成一个新数。

\textbf{样例入口。} 拿 [4, 1, 8, 7] 手算，先选 8 和 4 得到 4，再选 7 和 1 得到 6，最后 4*6 得到 24。若某一步选错，回溯要回到选两个数和运算符之前的状态继续试。

\textbf{暴力基线。} 慢版本就是完整搜索树：每层列出选择列表，路径到达终止条件时检查答案。它先保证覆盖所有可能，再把明显无效或重复的分支剪掉。放到本题就是：先按“四个数字必须全部使用，括号表达式可以看成每次选两个数合并成一个新数”完整试慢版本；样例里已经能看到“规律是任意括号表达式都能看成表达式树，最内层先合并两个叶子；状态是一组当前数字，每次合并两个数让集合大小减一”，所以优化前必须先能说清慢版本枚举了哪些候选。

\textbf{暴力过程。} 拿 [4, 1, 8, 7] 手算，先选 8 和 4 得到 4，再选 7 和 1 得到 6，最后 4*6 得到 24。若某一步选错，回溯要回到选两个数和运算符之前的状态继续试。

\textbf{重复工作。} 题面模型：四个数字必须全部使用，括号表达式可以看成每次选两个数合并成一个新数。暴力版的慢点要落到本题：慢版本会围绕“四个数字必须全部使用，括号表达式可以看成每次选两个数合并成一个新数”展开完整搜索树。样例规律是“规律是任意括号表达式都能看成表达式树，最内层先合并两个叶子；状态是一组当前数字，每次合并两个数让集合大小减一”，说明一层递归必须对应一个明确决策，剪枝只能删除整棵不可能成功或重复的子树。后面的优化只消掉这类重复工作，不能改变答案集合。

\textbf{优化条件。} 本题的关键观察是：从当前数字集合中枚举两个数和六种有序/无序运算，集合大小每次减少一，最后检查是否接近 24。这句话必须能翻译成可维护状态；如果题目条件改变后观察不再成立，就不能继续套原来的写法。

\textbf{相邻状态。} 规律是任意括号表达式都能看成表达式树，最内层先合并两个叶子；状态是一组当前数字，每次合并两个数让集合大小减一。把这个特例继续拆开看：一层递归对应一个明确决策，路径保存已经做出的选择，选择列表保存仍可能扩展的分支。剪枝前必须能说明被剪掉的整棵子树没有合法答案，而不是只因为它看起来不优。

\textbf{状态复用。} 把完整搜索树加上合法性检查、剩余资源剪枝和同层去重；路径状态进入和退出要成对恢复。本题落点是：规律是任意括号表达式都能看成表达式树，最内层先合并两个叶子；状态是一组当前数字，每次合并两个数让集合大小减一。手算时只改被新输入影响的那一小部分，而不是回头重算完整候选。

\textbf{指针和字段。} 状态字段要能说清：路径、起点或使用标记、剩余目标、答案集合、剪枝条件；进入递归和退出递归必须成对恢复。手算时按这个协议跟踪：拿 [4, 1, 8, 7] 手算，先选 8 和 4 得到 4，再选 7 和 1 得到 6，最后 4*6 得到 24。若某一步选错，回溯要回到选两个数和运算符之前的状态继续试。然后按协议复查：手算时给每层标出选择列表和路径。剪枝发生前要指出被剪分支为什么已经不可能得到合法答案。

\textbf{代码落点。} 剪枝条件写在扩展前；同层去重不要误写成全局去重；保存答案时复制当前路径。关键代码行必须能逐句翻译成“旧状态怎样变成新状态”，否则就只是模板。

\textbf{正确性。} 证明从搜索树覆盖开始：每个合法答案有且只有一条路径，剪枝只删掉不可能成功的分支。

\textbf{边界实验。} 先用减法和除法顺序、除零、浮点误差、重复中间状态做反例检查。实验时覆盖减法和除法顺序、除零、浮点误差、重复中间状态，统计搜索节点数量和剪枝次数。重复元素题要打印同一层跳过哪些值，确认没有误删合法路径。

\textbf{复杂度变化。} 暴力搜索树的规模由输出或选择空间决定；剪枝不改变最坏指数级本质，但能删除不可能成功或重复的分支，复杂度要说明输出规模。

\textbf{迁移追问。} 如果题目改成“可用排序后的状态去重剪枝，但核心仍是表达式树搜索”，先判断原状态是否仍然覆盖新答案集合，再决定是微调边界、改 value 语义，还是换算法。

\noindent\textbf{LeetCode 698 划分为 k 个相等的子集。}

\textbf{题目说明。} LeetCode 698 划分为 k 个相等的子集 的题面入口要先还原成具体任务：每个元素要分配到 k 个桶中，所有桶目标和相同。

\textbf{样例入口。} 拿 [4, 3, 2, 3, 5, 2, 1]、k=4 手算，总和 20，每桶目标 5；先放 5 单独成桶，再尝试 4+1、3+2。

\textbf{暴力基线。} 慢版本就是完整搜索树：每层列出选择列表，路径到达终止条件时检查答案。它先保证覆盖所有可能，再把明显无效或重复的分支剪掉。放到本题就是：先按“每个元素要分配到 k 个桶中，所有桶目标和相同”完整试慢版本；样例里已经能看到“规律是回溯状态是桶剩余容量和已用元素，排序后先放大数能更早剪枝”，所以优化前必须先能说清慢版本枚举了哪些候选。

\textbf{暴力过程。} 拿 [4, 3, 2, 3, 5, 2, 1]、k=4 手算，总和 20，每桶目标 5；先放 5 单独成桶，再尝试 4+1、3+2。

\textbf{重复工作。} 题面模型：每个元素要分配到 k 个桶中，所有桶目标和相同。暴力版的慢点要落到本题：慢版本会围绕“每个元素要分配到 k 个桶中，所有桶目标和相同”展开完整搜索树。样例规律是“规律是回溯状态是桶剩余容量和已用元素，排序后先放大数能更早剪枝”，说明一层递归必须对应一个明确决策，剪枝只能删除整棵不可能成功或重复的子树。后面的优化只消掉这类重复工作，不能改变答案集合。

\textbf{优化条件。} 本题的关键观察是：先判断总和可整除并排序降序，再用桶剩余容量回溯放置元素。这句话必须能翻译成可维护状态；如果题目条件改变后观察不再成立，就不能继续套原来的写法。

\textbf{相邻状态。} 规律是回溯状态是桶剩余容量和已用元素，排序后先放大数能更早剪枝。把这个特例继续拆开看：一层递归对应一个明确决策，路径保存已经做出的选择，选择列表保存仍可能扩展的分支。剪枝前必须能说明被剪掉的整棵子树没有合法答案，而不是只因为它看起来不优。

\textbf{状态复用。} 把完整搜索树加上合法性检查、剩余资源剪枝和同层去重；路径状态进入和退出要成对恢复。本题落点是：规律是回溯状态是桶剩余容量和已用元素，排序后先放大数能更早剪枝。手算时只改被新输入影响的那一小部分，而不是回头重算完整候选。

\textbf{指针和字段。} 状态字段要能说清：路径、起点或使用标记、剩余目标、答案集合、剪枝条件；进入递归和退出递归必须成对恢复。手算时按这个协议跟踪：拿 [4, 3, 2, 3, 5, 2, 1]、k=4 手算，总和 20，每桶目标 5；先放 5 单独成桶，再尝试 4+1、3+2。然后按协议复查：手算时给每层标出选择列表和路径。剪枝发生前要指出被剪分支为什么已经不可能得到合法答案。

\textbf{代码落点。} 剪枝条件写在扩展前；同层去重不要误写成全局去重；保存答案时复制当前路径。关键代码行必须能逐句翻译成“旧状态怎样变成新状态”，否则就只是模板。

\textbf{正确性。} 证明从搜索树覆盖开始：每个合法答案有且只有一条路径，剪枝只删掉不可能成功的分支。

\textbf{边界实验。} 先用存在元素大于目标、重复元素、空桶对称、k 为一做反例检查。实验时覆盖存在元素大于目标、重复元素、空桶对称、k 为一，统计搜索节点数量和剪枝次数。重复元素题要打印同一层跳过哪些值，确认没有误删合法路径。

\textbf{复杂度变化。} 暴力搜索树的规模由输出或选择空间决定；剪枝不改变最坏指数级本质，但能删除不可能成功或重复的分支，复杂度要说明输出规模。

\textbf{迁移追问。} 如果题目改成“状态压缩 DP 也能做，适合用掩码表示已选元素集合”，先判断原状态是否仍然覆盖新答案集合，再决定是微调边界、改 value 语义，还是换算法。

\subsection{实验复盘}

追问通常会问去重发生在同层还是路径、剪枝是否会删掉合法解、输出规模是多少。请准备画出前三层搜索树。

复盘本节时先从 LeetCode 17 电话号码的字母组合、LeetCode 22 括号生成、LeetCode 39 组合总和 开始：每道题都先手写一个能覆盖全部候选的暴力版，再指出优化结构具体保存了哪一段历史信息，最后用相邻案例比较为什么不能互换解法。
